\chapter{代数与方程}\label{ch:g9:algebra}

字母代表数，用字母计算可以证明关于
\emph{所有}数同时成立的事实。本章练习\omterm{def:g9:algebra:expand}{展开}与
因式分解——包括贯穿整个数学的三个恒等式——
并求解一次方程、积式方程与简单不等式，
每一步都有依据。

\section{展开}

\begin{definition}[展开]\label{def:g9:algebra:expand}
\emph{展开}\index{展开}是指利用分配律，将乘积化为
和的形式：
\[
k(a + b) = ka + kb,
\qquad
(a + b)(c + d) = ac + ad + bc + bd .
\]
\end{definition}

\begin{example}\label{ex:g9:algebra:expand}
第一个括号中的每一项都与第二个括号中的每一项相乘：
\begin{align*}
(x + 3)(2x + 5)
&= x \times 2x + x \times 5 + 3 \times 2x + 3 \times 5 \\
&= 2x^2 + 5x + 6x + 15 \\
&= 2x^2 + 11x + 15 .
\end{align*}
符号随各项一同移动：
\begin{align*}
(3x - 2)(x - 4)
&= 3x^2 - 12x - 2x + 8
= 3x^2 - 14x + 8 .
\end{align*}
\end{example}

\begin{theorem}[三个恒等式]\label{thm:g9:algebra:identities}
对所有数 $a$ 与 $b$：
\[
(a+b)^2 = a^2 + 2ab + b^2,
\qquad
(a-b)^2 = a^2 - 2ab + b^2,
\qquad
(a+b)(a-b) = a^2 - b^2 .
\]
\end{theorem}

\begin{proof}
分别\omterm{def:g9:algebra:expand}{展开}每个等式的左边。对第一式，
$(a+b)(a+b) = a^2 + ab + ba + b^2 = a^2 + 2ab + b^2$。对第二式，
将第一式中的 $b$ 换为 $-b$。对第三式：
$(a+b)(a-b) = a^2 - ab + ba - b^2 = a^2 - b^2$（交叉项相消）。
\end{proof}

\begin{omfigure}
\begin{tikzpicture}[scale=0.60]
  \draw[thick] (0,0) rectangle (5,5);
  \draw[thick] (3.3,0) -- (3.3,5);
  \draw[thick] (0,3.3) -- (5,3.3);
  \fill[omDef!14] (0,0) rectangle (3.3,3.3);
  \fill[omThm!14] (3.3,3.3) rectangle (5,5);
  \fill[omProp!14] (3.3,0) rectangle (5,3.3);
  \fill[omProp!14] (0,3.3) rectangle (3.3,5);
  \node[font=\small] at (1.65,1.65) {$a^2$};
  \node[font=\small] at (4.15,4.15) {$b^2$};
  \node[font=\small] at (4.15,1.65) {$ab$};
  \node[font=\small] at (1.65,4.15) {$ab$};
  \node[below, font=\small] at (1.65,0) {$a$};
  \node[below, font=\small] at (4.15,0) {$b$};
  \node[left, font=\small] at (0,1.65) {$a$};
  \node[left, font=\small] at (0,4.15) {$b$};
\end{tikzpicture}

{\small 恒等式 $(a+b)^2 = a^2 + 2ab + b^2$，从\omterm{def:g6:measure:area}{面积}看：边长为 $a+b$ 的大
正方形由两个正方形和两个相等的长方形组成。
忘记“$2ab$”，就等于忘了那两个长方形！}
\end{omfigure}

\begin{example}\label{ex:g9:algebra:identities}
\[
(x + 5)^2 = x^2 + 10x + 25, \qquad
(3x - 1)^2 = 9x^2 - 6x + 1, \qquad
(2x + 7)(2x - 7) = 4x^2 - 49 .
\]
口算：$102^2 = (100 + 2)^2 = 10000 + 400 + 4 = 10404$，又
$98 \times 102 = (100-2)(100+2) = 10000 - 4 = 9996$。
\end{example}

\section{因式分解}

\begin{method}[因式分解]\label{met:g9:algebra:factor}
因式分解一个式子时，依次尝试：
\begin{enumerate}
  \item \emph{公因式}：$ka + kb = k(a + b)$，其中 $k$ 本身
        可以是一个括号；
  \item \emph{平方差}：$a^2 - b^2 = (a+b)(a-b)$；
  \item \emph{完全平方}：$a^2 + 2ab + b^2 = (a+b)^2$（若中间项为
        $-2ab$，则为 $(a-b)^2$）。
\end{enumerate}
总要把结果\omterm{def:g9:algebra:expand}{展开}来检验。
\end{method}

\begin{example}\label{ex:g9:algebra:factor}
\emph{公因式：}
$15x^2 + 10x = 5x \times 3x + 5x \times 2 = 5x(3x + 2)$。

\emph{公共括号：}
\[
(x+2)(x-1) + (x+2)(3x+4) = (x+2)\bigl[(x-1) + (3x+4)\bigr]
= (x+2)(4x+3).
\]

\emph{平方差：}
$x^2 - 81 = (x+9)(x-9)$，且
\[
(2x+1)^2 - 9 = (2x+1)^2 - 3^2 = (2x+1+3)(2x+1-3) = (2x+4)(2x-2).
\]
各因式中仍有公因式：$(2x+4)(2x-2) = 2(x+2) \times
2(x-1) = 4(x+2)(x-1)$。
\end{example}

\section{方程}

\begin{method}[一元一次方程]\label{met:g9:algebra:firstdegree}
求解形如 $5x - 3 = 2x + 9$ 的方程：
\begin{enumerate}
  \item 将含 $x$ 的项移到方程一边，常数项移到另一边，
        方法是对两边加上同一个量；
  \item 化简方程两边；
  \item 方程两边同除以 $x$ 的系数；
  \item 将所得解代入原方程进行检验。
\end{enumerate}
\end{method}

\begin{example}\label{ex:g9:algebra:firstdegree}
\begin{align*}
5x - 3 &= 2x + 9 \\
5x - 2x &= 9 + 3 && \text{（两边同时加上 $3 - 2x$）}\\
3x &= 12 \\
x &= 4 && \text{（两边同除以 3）。}
\end{align*}
检验：$5 \times 4 - 3 = 17$，$2 \times 4 + 9 = 17$。方程的解为
$4$。

含括号时，先\omterm{def:g9:algebra:expand}{展开}：$3(x - 2) = x + 8$ 化为 $3x - 6 = x + 8$，
于是 $2x = 14$，$x = 7$。
\end{example}

\begin{theorem}[零因子法则]\label{thm:g9:algebra:zeroproduct}
两个数的积为零，当且仅当至少有一个因式为零：
\[
A \times B = 0
\quad\Longleftrightarrow\quad
A = 0 \ \text{ 或 } \ B = 0 .
\]
\end{theorem}

\begin{proof}
若某一因式为零，则积显然为零。反之，若 $A B = 0$ 且
$A \neq 0$，两边同除以 $A$ 得 $B = 0$。
\end{proof}

\begin{example}\label{ex:g9:algebra:zeroproduct}
解方程 $(x - 3)(2x + 10) = 0$：或 $x - 3 = 0$，或 $2x + 10 = 0$，故
解为 $3$ 与 $-5$。

解方程 $x^2 - 49 = 0$：先因式分解，$(x-7)(x+7) = 0$，解为 $7$ 与
$-7$。

解方程 $x^2 = 3x$：切勿两边同除以 $x$！移项并分解：
$x^2 - 3x = x(x - 3) = 0$，解为 $0$ 与 $3$。
\end{example}

\section{不等式}

\begin{proposition}[不等式的性质]\label{prop:g9:algebra:ineq}
不等式两边同时加上同一个数时，不等号方向不变，
两边同乘同一个\emph{正数}时，同样成立。
两边同乘一个\emph{负数}时，
不等号方向改变。
\end{proposition}

\begin{proof}
若 $a \leq b$，则 $b - a \geq 0$。加上 $c$：$(b+c) - (a+c) = b - a
\geq 0$，故 $a + c \leq b + c$。乘以 $c > 0$：
$bc - ac = (b-a)c \geq 0$，故 $ac \leq bc$。乘以 $c < 0$：
$(b-a)c \leq 0$，故 $ac \geq bc$，不等号方向改变。
\end{proof}

\begin{example}\label{ex:g9:algebra:ineq}
解不等式 $7 - 2x < 15$：
\begin{align*}
7 - 2x &< 15 \\
-2x &< 8 && \text{（两边同减 7）} \\
x &> -4 && \text{（两边同除以 $-2$，不等号方向改变）。}
\end{align*}
解集为所有大于 $-4$ 的数：在数轴上，
以空心点表示 $-4$，并向右涂色。
\end{example}

\begin{omfigure}
\begin{tikzpicture}[scale=1.0]
  \draw[-Stealth] (-6.4,0) -- (2.6,0) node[below, font=\small] {$x$};
  \foreach \x in {-6,-5,-4,-3,-2,-1,0,1,2}
    \draw (\x,-0.07) -- (\x,0.07);
  \node[below, font=\small] at (-4,-0.12) {$-4$};
  \node[below, font=\small] at (0,-0.12) {$0$};
  \draw[omDef, very thick] (-4,0.02) -- (2.5,0.02);
  \draw[omDef, fill=white, thick] (-4,0.02) circle (2.4pt);
\end{tikzpicture}

{\small 不等式 $7 - 2x < 15$ 的解：所有满足 $x > -4$ 的数（空心点：$-4$
本身不是解）。}
\end{omfigure}

\section{练习}

\begin{exercise}[$\star$]\label{exo:g9:algebra:1}
\omterm{def:g9:algebra:expand}{展开}并化简：
\[
4(3x - 2), \qquad
(x + 6)(x + 2), \qquad
(2x - 3)(x + 5), \qquad
(x - 1)(x - 9).
\]
\end{exercise}

\begin{exercise}[$\star$]\label{exo:g9:algebra:2}
利用恒等式\omterm{def:g9:algebra:expand}{展开}：
\[
(x + 8)^2, \qquad
(3x - 4)^2, \qquad
(x + 10)(x - 10), \qquad
(5x + 2)(5x - 2).
\]
\end{exercise}

\begin{exercise}[$\star$]\label{exo:g9:algebra:3}
利用恒等式心算：$101^2$；\quad $59^2$
（写成 $59 = 60 - 1$）；\quad $53 \times 47$。
\end{exercise}

\begin{exercise}[$\star$]\label{exo:g9:algebra:4}
因式分解：
\[
9x + 12, \qquad
x^2 - 64, \qquad
x^2 + 14x + 49, \qquad
18x^2 - 6x .
\]
\end{exercise}

\begin{exercise}[$\star$]\label{exo:g9:algebra:5}
解下列方程，写出每一步：
\[
4x + 7 = 19, \qquad
6x - 5 = 2x + 11, \qquad
5(x - 3) = 3x + 1 .
\]
\end{exercise}

\begin{exercise}[$\star\star$]\label{exo:g9:algebra:6}
利用零乘积性质求解：
\[
(x + 4)(3x - 12) = 0, \qquad
x^2 - 121 = 0, \qquad
(2x - 1)^2 = 0, \qquad
x^2 = 8x .
\]
\end{exercise}

\begin{exercise}[$\star\star$]\label{exo:g9:algebra:7}
因式分解 $E = (3x + 2)^2 - 25$，再解方程 $E = 0$。
\end{exercise}

\begin{exercise}[$\star\star$]\label{exo:g9:algebra:8}
解下列不等式，并在数轴上表示解集：
\[
3x - 4 \geq 8, \qquad
5 - 2x > 1, \qquad
4(x + 1) \leq x - 5 .
\]
\end{exercise}

\begin{exercise}[$\star\star$]\label{exo:g9:algebra:9}
想一个数，把它乘以 $3$，再加上 $7$，所得结果与
将它乘以 $5$ 再减去 $9$ 所得的结果一样。这个数是多少？
（列方程并求解。）
\end{exercise}

\begin{exercise}[$\star\star\star$]\label{exo:g9:algebra:10}
设 $n$ 为整数。
\begin{enumerate}
  \item \omterm{def:g9:algebra:expand}{展开} $(n + 1)(n - 1)$，并由此推出 $n$ 的两个邻数之积
        总是 $n^2 - 1$（例如 $24 \times 26 =
        25^2 - 1 = 624$）。
  \item 证明：三个连续整数之和总是
        $3$ 的\omterm{def:g9:arith:divisor}{倍数}。
\end{enumerate}
\end{exercise}

\section{问题：数字的代数}

\begin{problem}[{周末问题——整除规则为何成立，
1089 戏法的揭秘，以及
看似魔法的心算}]
\label{pb:g9:algebra:1}
小学里你学过\emph{规则}（一个数能被
$9$ \omterm{def:g9:arith:divisor}{整除}，当且仅当其数字和亦然）；魔术师表演的是\emph{戏法}
（人人最后都得到 $1089$）；市场上的人懂\emph{捷算}
（$35^2 = 1225$，瞬间即得）。它们都是同一个秘密：一个
数字依次为 $a$、$b$、$c$ 的数，可写成表达式
$100a + 10b + c$，而本章的代数
（\cref{thm:g9:algebra:identities}，\omterm{def:g9:algebra:expand}{展开}与因式分解）可以
把它拆开。学完本题后，你将证明
这些规则，揭穿这一戏法，并学会其中的“魔法”。

\medskip
\noindent\textbf{第一部分——显微镜下的数字。}
\begin{enumerate}
  \item 设十位数字为 $a$、个位数字为 $b$ 的两位数
        是 $10a + b$。\omterm{def:g9:algebra:expand}{展开}并因式分解这个数与其反序数
        之差，
        $(10a + b) - (10b + a)$，并陈述你的发现。以
        $72 - 27$ 为例进行检验。
  \item 证明三位数的九的判别法：验证
        恒等式
        \[
        100a + 10b + c = 9\,(11a + b) + (a + b + c),
        \]
        并说明为何这表明一个数与其数字和
        在除以 $9$ 时留下\emph{相同的\omterm{def:g3:division:remainder}{余数}}
——因此一个能被 $9$ \omterm{def:g9:arith:divisor}{整除}，当且仅当另一个亦然
        （\cref{prop:g9:arith:rules}，现已证明）。
  \item 由同一恒等式推出三的判别法，
        用一句话说明。
  \item “弃九法”，即会计师的验算方法：检验
        $47 \times 86 = 4\,042$，将每个数替换为
        其数字和的数字和：$47 \to 11 \to 2$，
        $86 \to 14 \to 5$，$2 \times 5 = 10 \to 1$；然后
        $4\,042 \to 10 \to 1$：结果一致。说明为何正确的
        积必定通过检验（设 $47 = 9k + 2$，
        $86 = 9l + 5$，并\omterm{def:g9:algebra:expand}{展开}），然后找出一个
        检验无法发现的错误结果（为何交换数字对它
        是不可见的？）。
  \item 证明四位数的十一的判别法：验证
        \[
        1000a + 100b + 10c + d
        = 11\,(91a + 9b + c) + (- a + b - c + d),
        \]
        并用它判断 $2\,926$ 是否能被
        $11$ \omterm{def:g9:arith:divisor}{整除}。
\end{enumerate}

\medskip
\noindent\textbf{第二部分——1089 魔术。}
魔术：想一个三位数，它的首位与末位数字
至少相差 $2$；把这个数倒过来写；用较大的数减去较小的数；
把所得的差倒过来写；把最后两个数相加。
魔术师宣布：\emph{1089}。
\begin{enumerate}[resume]
  \item 对 $742$ 和你自选的一个数，
        分别完整地做一遍该魔术。
  \item 设该数的三个数字依次为 $a$、$b$、$c$，且 $a > c$。
        将第一次减法因式分解，
        $(100a + 10b + c) - (100c + 10b + a)$，并列出
        当 $a - c$ 取 $2$ 至 $9$ 时它所能取的全部值。
  \item 这些 $99$ 的\omterm{def:g9:arith:divisor}{倍数}都具有醒目的
        数字规律。记 $m = a - c$，验证恒等式
        \[
        99m = 100\,(m - 1) + 90 + (10 - m),
        \]
        并读出 $99m$ 的三位数字。
  \item 把 $99m$ 与它倒过来的数相加，
        利用第 8 题中的数字，证明其和为 $1089$——
        无论 $m$ 为何值。魔术已不再神秘；代数万岁！
  \item 魔术师为何要求首位与末位数字
        至少相差 $2$？考察情形 $a = c$ 与
        $a - c = 1$（此时减法结果为 $99$，
        须写成 $099$，倒序才能成立），并
        写出保证魔术可靠的附加条件。
\end{enumerate}

\medskip
\noindent\textbf{第三部分——日常生活中的魔术。}
\begin{enumerate}[resume]
  \item 生日魔术：“把你的出生月份乘以 $5$，
        加 $7$，翻倍，再加上你的出生日，最后减去 $14$。”
        证明结果恒为 $10m + d$——即月份与
        日期，一目了然。（可以找人试一试。）
  \item 你将该魔术改为：“把月份乘以 $5$，
        加 $9$，翻倍，再加上日期。”此时观众
        最后必须减去哪个数，才能得到相同的
        月日读数？
  \item 缺失数字魔术：观众任选一个数，
        减去它的数字和（第 2 题指出该结果
        总是 $9$ 的\omterm{def:g9:arith:divisor}{倍数}——请对一个
        三位数解释其原因），然后划去一个\emph{非零}
        数字，并把剩下的数字告诉你。说明如何
        找回被划去的数字，
        以及为何必须排除 $0$。
  \item 市场魔术：证明恒等式
        $(10t + 5)^2 = 100\,t(t+1) + 25$，叙述
        任一个以 $5$ 结尾的数的平方法则，并心算 $45^2$、
        $85^2$ 与 $105^2$。
  \item 压轴题，推广\cref{exo:g9:algebra:10}：证明
        任意五个连续整数之和总是
        $5$ 的\omterm{def:g9:arith:divisor}{倍数}，但\emph{四}个连续整数之和
        \emph{从来不}是 $4$ 的\omterm{def:g9:arith:divisor}{倍数}。
        （两段简短的计算——以及一条关于代数
        代数如何处理“总是”与“从来不”的一条道理。）
\end{enumerate}
\end{problem}
